\subsection{The finite-weight collision bound in rank four}
\label{sec:rh-rank-four-weight-collision}

We isolate the finite-weight part of the component-propagation argument.
Unlike the construction of the ordinary automorphic ideal, this part
requires no automorphy lifting. The source statement is
\texttt{off\_diagonal\_nilpotent\_rank\_four} in
\texttt{OffDiagonalNilpotentRankFour.lean}.

\begin{theorem}[Dimension loss at a collision of inertial weights]
\label{thm:rh-rank-four-weight-collision}
Retain the original fitted deep tower \(T\), odd prime \(\ell\),
coefficient DVR \(\Lambda=c'.O\) with uniformizer \(\varpi\),
and original normalized space \(X\) with coefficient ring \(B=X.B\).
Write \(d=[T.F:\mathbb Q]\), \(r=\#T.\iota\), and
\(d_v=[T.F_v:\mathbb Q_\ell]\) for the common local degree at the
\(r\) selected places, so \(rd_v=d\).
Let \(\psi_{u,i}(s)=X.\psi(u,i,s)\) be the original diagonal
local characters, and \(J\subset B\) an ideal satisfying the original
finite-weight property
\[
 \mathfrak m_B^n\subset
 J+\varpi B+
 (\psi_{u,i}(s)-1:u\in T.\iota,\ 0\leq i\leq3,\ s\in I_u)
\]
for some \(n\), where \(I_u\) denotes the inertia group at \(u\).
No other part of \texttt{PotentialModularityTower} is needed.

For distinct \(i,j\), a place \(v\), and a prime
\(\mathfrak q\supset J+\varpi B\) such that
\[
 \psi_{v,i}(s)-\psi_{v,j}(s)\in\mathfrak q
 \qquad(s\in T.U1(v)),
\]
one has the literal source bound
\[
              \dim B/\mathfrak q+d_v\leq2d+1 .
\]
\end{theorem}

\begin{proof}
Write \(A=B/\mathfrak q\). We use the following elementary facts on
all further prime quotients of \(A\). Such a quotient \(S\) is a
complete Noetherian local domain of characteristic \(\ell\), with
finite residue field. By definition of the original space, the
images of \(\psi_{u,i}\) are adically continuous characters on the
decomposition groups. Indeed, the local upper triangular matrix
equations give multiplicativity of the diagonal entries; their
products in opposite positions equal the prescribed similitude.
These identities descend from \(X.A\) to \(B\) by the injective
coefficient map, and continuity follows from the original
\texttt{psi\_cont} condition. The diagonal residual characters on
inertia are all trivial by \(T.Ok.\mathrm{ineL}\), so the inertia
characters \(\chi_{u,i}\) of \(S\) take values in \(1+\mathfrak m_S\)
and satisfy
\[
                   \chi_{u,i}\chi_{u,3-i}=1 .
\]
Here the multiplier is one on inertia after reducing the coefficient
DVR, again by the same \(\mathrm{ineL}\) condition. These arguments
give the relevant part of
\path{deformationSpace_local_characters_continuous} directly on
the original space; no change to its deformation functor is made.

The group \(1+\mathfrak m_S\) is pro-\(\ell\): the successive
congruence quotients are additive groups of finite characteristic
\(\ell\). It is torsion-free. For \(\ell\)-power torsion this follows
from \((z-1)^{\ell^a}=z^{\ell^a}-1\) and the domain condition; for
prime-to-\(\ell\) torsion, if \(z\equiv1\bmod\mathfrak m_S\) and
\(z^b=1\), \(b\) prime to \(\ell\), then
\(1+z+\cdots+z^{b-1}\) is a unit and \(z=1\).

Local class field theory identifies the image of inertia in the
abelianized decomposition group with the unit group of \(T.F_u\).
Its maximal pro-\(\ell\) quotient modulo finite torsion is free
of rank \(d_v\) over \(\mathbb Z_\ell\). Consequently every
\(\chi_{u,i}|_{I_u}\) factors continuously through this same
torsion-free quotient \(L_u\). Choose inertia lifts
\(s_{u,1},\dots,s_{u,d_v}\) of a basis of \(L_u\). Vanishing of
\(\psi_{u,i}(s_{u,k})-1\) in \(S\) for all \(k\) implies that
\(\chi_{u,i}\) is trivial on inertia. The lifts depend only on
the tower, not on the prime quotient.

At \(v\), the cyclotomic character on the pro-\(\ell\) part of the
unit group is nontrivial and has open image in
\(1+\ell\mathbb Z_\ell\); for example this is the local norm on
units up to the reciprocity convention, and a sufficiently deep
principal-unit subgroup has open norm. It kills finite torsion.
The induced map on \(L_v\) has kernel a direct summand of
rank \(d_v-1\) and quotient free of rank one. The image of
\(T.U1(v)\) contains that kernel: a lift in the principal-unit
group of a kernel element has trivial cyclotomic value,
and every relative inertia element has trivial original residual
representation by \(\mathrm{ineL}\).
This uses precisely the source definition
\[
 T.U1(v)=\{s\in I_v:\epsilon_\ell(s)=1,\
                         D.\bar r(\operatorname{res}_{\mathbb Q}s)=1\}.
\]
Choose \(s_*\in I_v\) in the principal-unit part whose image
generates the free rank-one quotient. It follows that the
collision in the theorem, together with the single relation
\(\psi_{v,i}(s_*)-\psi_{v,j}(s_*)=0\), forces
\(\chi_{v,i}=\chi_{v,j}\) on all inertia in every prime quotient
of \(A\). This explains the one additional parameter in the source
bound. Torsion units introduce no additional parameters.

There are two opposite pairs of labels, \(\{0,3\}\) and
\(\{1,2\}\). Choose a label \(a\) in the pair other than the pair
of \(i\), taking \(a=j\) if \(j\ne3-i\). Consider the following
elements of \(A\):
\[
\begin{array}{ll}
 \psi_{u,a}(s_{u,k})-1
       & (u\in T.\iota,\ 1\leq k\leq d_v),\\
 \psi_{u,i}(s_{u,k})-1
       & (u\ne v,\ 1\leq k\leq d_v),\\
 \psi_{v,i}(s_*)-\psi_{v,j}(s_*).&
\end{array}
\]
All are in the maximal ideal. Their number is
\(rd_v+(r-1)d_v+1=2d-d_v+1\).
Let \(\mathfrak p\) be any prime of \(A\) containing them and set
\(S=A/\mathfrak p\). The first family forces the character of
\(a\), and hence of \(3-a\), to be trivial on inertia at every
place. The second family does the same for \(i,3-i\) away from
\(v\). At \(v\) the third relation completes the inertial
collision. If \(j=a\), it makes the character of \(i\) trivial;
if \(j=3-i\), it forces \(\chi_{v,i}^2=1\), which also makes it
trivial since the target \(1+\mathfrak m_S\) is torsion-free.
The opposite character is then trivial as well.

Thus every inertial parameter in the finite-weight hypothesis
vanishes in \(S\). Since \(J,\varpi B\) already vanish, that
hypothesis implies \(\mathfrak m_S^n=0\); the domain \(S\) is a
field. Hence the displayed \(2d-d_v+1\) parameters generate
an ideal of \(A\) whose radical is its maximal ideal. Krull's
height theorem gives \(\dim A\leq2d-d_v+1\), as asserted.
\end{proof}

For the actual component step, the source chooses
\(N=2d+1-2r\) with \(r\leq C\) and
\(30(C+1)^2\leq d\). In this range \(d_v=d/r\geq30(C+1)\)
and \(N\geq d+20C+2\), so the strict inequalities needed by
\texttt{ch\_core} follow:
\[
 d+20C+1<N,\qquad
 2d+1+(2+4r)<N+d_v.
\]
The second is equivalent to \(d_v>2+6r\).
These checks retain the original tower threshold.
Theorem~\ref{thm:rh-rank-four-weight-collision} supplies its
weight-collision estimate. The characteristic-prime curve
propagation still needs the actual local chart and reducible-locus
and irreducible-fibre dimension bounds, and the construction of one
ideal \(J\) with the remaining automorphic properties.
In particular, we have not proved \texttt{PMExAlt} or
\texttt{residual\_charpoly\_rank\_four} by this argument.
